\part{Reference}

\chapter{Préréglages de confort et visibilité}
\label{ch:comfort-presets}

\rowcolors{2}{LtGray}{white}
\begin{longtable}{C{1.4cm}L{2.0cm}L{4.0cm}L{6.7cm}}
\toprule
\rowcolor{Navy}
\color{white}\sffamily\bfseries Icon &
\color{white}\sffamily\bfseries Preset &
\color{white}\sffamily\bfseries Conditions &
\color{white}\sffamily\bfseries Key characteristics \\
\midrule
\iconlg{comfort-default} & Default & General fallback &
  Balanced mid-contrast; usable in most conditions. \\
\iconlg{comfort-dawn}    & Dawn    & Around sunrise &
  Warm tones, slightly reduced brightness. \\
\iconlg{comfort-day}     & Day     & Full daylight, direct sun &
  Maximum brightness and contrast. For open-ocean sailing. \\
\iconlg{comfort-sunset}  & Sunset  & Around sunset &
  Warmer, transitional brightness. \\
\iconlg{comfort-dusk}    & Dusk    & Twilight &
  Reduced brightness; begins night-vision preservation. \\
\iconlg{comfort-night}   & Night   & Full dark, offshore watches &
Brillance minimale, fond sombre, icônes teintées rouge.
Préserve la vision nocturne.
  Alarms remain high-contrast on every preset. \\
\bottomrule
\end{longtable}

\section{Orientations pratiques}

\begin{itemize}
  \item Use brighter, higher-contrast presets in daylight.
  \item Use dimmer, night-oriented presets when preserving night vision matters.
  \item Keep alarms and safety-critical text easy to read in every preset.
  \item Prefer readability over appearance. If a preset looks attractive but makes
information urgente plus difficile à voir, passer à une plus claire.
  \item The Night preset sets the icon tint to red by default.
C'est intentionnel --- la lumière rouge minimise le temps nécessaire pour
ré-adapter à l'obscurité après avoir regardé l'affichage.
\end{itemize}

\chapter{Opération d'affichage, de contact et de conditions difficiles}
\label{ch:roughconditions}

\section{Afficher la lisibilité}

FairWindSK devrait rester lisible en plein jour, crépuscule, nuit et rugueux
conditions.
Si l'affichage est difficile à lire, corriger d'abord avant de diagnostiquer n'importe quel
problèmes d'application.

Commencez par :
\begin{enumerate}
  \item Screen brightness.
  \item Comfort preset (correct for the current light conditions).
  \item Viewing angle and helm position.
  \item The app or page shown in the Application Area.
\end{enumerate}

Gardez les numéros critiques de sécurité faciles à identifier en un coup d'oeil.
Évitez les préréglages qui réduisent le contraste trop.
Utilisez des écrans plus grands pour la navigation primaire dans la mesure du possible.
Utilisez un parasol ou une hotte en plein soleil.

\section{Utilisation du toucher dans des conditions difficiles}

\begin{itemize}
  \item Brace yourself first. One hand on the boat, one on the screen.
  \item Use deliberate, single taps. Hesitant pressure and repeated rapid taps
provoquer des actions involontaires.
  \item Never explore the autopilot controls by tapping speculatively ---
chaque commande a un effet physique immédiat.
  \item Keep the screen clean. Salt spray reduces touch accuracy significantly.
  \item Pre-configure what you need while conditions are calm.
  \item Prefer the large, stable targets of the Bottom Bar main icons.
  \item Confirm dialogs and mode changes before acting on them.
\end{itemize}

Pour l'utilisation de la barre, la disposition stable et une bonne lisibilité sont plus importantes que le style décoratif.

\chapter{Flux de travail recommandé de l'opérateur}
\label{ch:workflow}

Pour la plupart des passages, cette séquence maintient le shell prévisible et facile à faire confiance.

\section{Avant le départ}

\begin{enumerate}
  \item Start FairWindSK and confirm the shell health state is normal.
  \item Verify that expected key data is live in the Top Bar.
  \item Open the app most relevant to the current passage phase.
  \item Confirm brightness and comfort preset.
  \item Complete the pre-departure checklist (Chapter~2).
\end{enumerate}

\section{Pendant le passage}

\begin{enumerate}
  \item Watch for stale or missing values in the Top Bar.
  \item Cross-check route, depth, and heading whenever conditions change.
  \item Use the launcher to switch tasks rather than forcing a stuck app.
  \item Treat degraded states as operational warnings.
  \item If the current screen stops feeling clear or trustworthy, tap
        \iconlg{applications}~\key{Apps} to return to the launcher.
\end{enumerate}

\section{Après une reconnexion ou un redémarrage du serveur}

\begin{enumerate}
  \item Wait briefly for automatic recovery to complete.
  \item Confirm the status shows \textit{Live}.
  \item Reopen the app if the foreground application was degraded.
  \item Reconfirm key navigation values before depending on them.
\end{enumerate}

\section{Changement de tâche}

Utilisez la barre inférieure plutôt que de forcer l'application actuelle à tout faire :

\begin{itemize}
  \item Use \iconlg{applications}~Apps to return to the launcher and switch tasks.
  \item Use the open-apps strip to jump directly to a loaded application.
  \item Use operational bars (\iconlg{alarm-pob}~POB, \iconlg{anchor-iec}~Anchor,
        \iconlg{command-autopilot}~Autopilot) for vessel-operation tasks
sans quitter votre application de navigation.
\end{itemize}

\chapter{Serveur Signal K sur Docker}
\label{ch:docker}

\section{Image officielle Docker}

\begin{lstlisting}
cr.signalk.io/signalk/signalk-server:latest
\end{lstlisting}

Supported architectures: \code{linux/amd64}, \code{linux/arm/v7},
\code{linux/arm64} (all Raspberry~Pi~4 and~5 configurations).

\section{Démarrer rapidement avec les données de la démonstration NMEA}

\begin{lstlisting}
mkdir $HOME/.signalk
docker run --init -it --rm \
  --name signalk-server \
  --publish 3000:3000 \
  -v $HOME/.signalk:/home/node/.signalk \
  cr.signalk.io/signalk/signalk-server \
  --sample-nmea0183-data
\end{lstlisting}

\section{Déploiement de production}

\begin{lstlisting}
docker run -d --init \
  --name signalk-server \
  -p 3000:3000 \
  -v $HOME/.signalk:/home/node/.signalk \
  cr.signalk.io/signalk/signalk-server
\end{lstlisting}

\section{Mise à jour du signal K}

\begin{lstlisting}
docker pull cr.signalk.io/signalk/signalk-server:latest
docker stop signalk-server
docker rm signalk-server
# Re-run the production docker run command above
\end{lstlisting}

\chapter{Dépannage}
\label{ch:troubleshooting}

\section{Je vois des valeurs manquantes}

Causes possibles : le capteur ou l'instrument source est hors ligne ; Signal~K n'est pas
réception de ce chemin; la valeur n'est pas configurée pour l'affichage; la connexion a
tombé.

\begin{enumerate}
  \item Check shell health state.
  \item See whether other values are still live.
  \item Confirm the source instrument is operating.
  \item Check \ui{Settings~> Signal~K Paths} --- is the correct path configured?
  \item Return to the launcher and reopen the affected app if needed.
\end{enumerate}

\section{Je vois les valeurs de stales}

Causes possibles: interruption temporaire des données; reconnectation en cours; retard du réseau;
dégradation partielle de la source de données.

\begin{enumerate}
  \item Treat the values cautiously.
  \item Wait briefly for recovery.
  \item Cross-check with other instruments.
  \item Investigate if the stale state persists beyond 60 seconds.
\end{enumerate}

\section{Une application Web est vierge ou congelée}

\begin{enumerate}
  \item Retry once if the shell offers the option.
  \item Return to the launcher.
  \item Check shell health and Signal~K connection state.
  \item Use another app if the shell is otherwise healthy.
  \item A failed application does not mean FairWindSK has failed.
\end{enumerate}

\section{Le Shell est déconnecté}

\begin{enumerate}
  \item Is Signal~K running? Check LEDs on the Raspberry~Pi; check the Docker
        container status: \code{docker ps -a | grep signalk}.
  \item Is your device on the correct Wi-Fi network?
  \item Is the server address in \ui{Settings~> Connection} correct?
  \item Can a browser on another device reach \code{http://server-address:3000}?
  \item Allow FairWindSK 60 seconds to attempt automatic reconnection.
  \item Reconfirm all safety-critical values after reconnection.
\end{enumerate}

\section{Aucune application n'apparaît sur le lanceur}

\begin{enumerate}
  \item Confirm status shows \textit{Live}.
  \item Are apps installed in the Signal~K App Store?
  \item In \ui{Settings~> Connection}, tap \key{Connect} to force a reconnect and
app découverte.
\end{enumerate}

\section{Le pilote automatique ou l'icône d'ancrage est grisé}

\begin{enumerate}
  \item Is the required plugin installed in the Signal~K App Store?
  \item Is the plugin enabled and configured in Signal~K settings?
  \item Is the correct plugin action path configured in
        \ui{Settings~> Signal~K Paths}?
\end{enumerate}

\section{FairWindSK roule lentement sur Raspberry Pi}

\begin{enumerate}
  \item Check CPU temperature: \code{vcgencmd measure\_temp}.
Au-dessus de 80°C indique un étranglement thermique --- améliorer le refroidissement.
  \item Check \ui{Settings~> System~> Performance} for memory usage.
        FairWindSK should use under 250\,MB at idle.
  \item Use a Class~A2 high-speed microSD card.
Une carte lente est la cause la plus courante de la performance lugubre sur Raspberry~Pi.
  \item Set Log Level to \textit{Off} in \ui{Settings~> System~> Diagnostics}
pour minimiser les écritures de carte SD.
\end{enumerate}

\chapter{Liste quotidienne de contrôle d'exploitation}
\label{ch:checklist}

\section*{Before Departure}

\rowcolors{2}{LtGray}{white}
\begin{longtable}{C{0.8cm}L{12.9cm}}
\toprule
\rowcolor{Navy}
\color{white}\sffamily\bfseries $\square$ &
\color{white}\sffamily\bfseries Check \\
\midrule
$\square$ & FairWindSK started --- status shows \textit{Live} \\
$\square$ & Position live and plausible \\
$\square$ & SOG updating --- valid GPS fix \\
$\square$ & Depth updating (if fitted) \\
$\square$ & Heading showing \\
$\square$ & Chart plotter open; vessel in correct position \\
$\square$ & Route or waypoint active (if proceeding to a destination) \\
$\square$ & AIS targets visible (if AIS fitted) \\
$\square$ & Comfort preset set for current light conditions \\
$\square$ & Screen brightness readable from helm \\
$\square$ & POB button location confirmed \\
$\square$ & Autopilot available and responding (if required) \\
$\square$ & Anchor watch available (if anchoring tonight) \\
\bottomrule
\end{longtable}

\section*{After Arrival}

\rowcolors{2}{LtGray}{white}
\begin{longtable}{C{0.8cm}L{12.9cm}}
\toprule
\rowcolor{Navy}
\color{white}\sffamily\bfseries $\square$ &
\color{white}\sffamily\bfseries Check \\
\midrule
$\square$ & Return to launcher --- no active alarms \\
$\square$ & Anchor watch set with appropriate radius (if anchoring) \\
$\square$ & Comfort preset changed to Dusk or Night \\
$\square$ & Track exported for logbook (if required) \\
\bottomrule
\end{longtable}

\chapter{Carte de référence rapide}
\label{ch:quickref}

\section*{Immediate Action Guide}

\rowcolors{2}{LtGray}{white}
\begin{longtable}{L{3.6cm}L{4.4cm}L{5.7cm}}
\toprule
\rowcolor{Navy}
\color{white}\sffamily\bfseries Situation &
\color{white}\sffamily\bfseries Immediate action &
\color{white}\sffamily\bfseries Then \\
\midrule
Personne hors pension&
  Tap \iconlg{alarm-pob}~\key{POB}. Once. &
  Call MAYDAY on VHF Ch~16. Deploy SAR route. \\
Alarme clignotante&
  Tap \iconlg{alarm-general}~\key{Alarms}. Read it. &
  Act per emergency procedures. \\
Écran gelé&
  Wait 10~s. Tap \iconlg{applications}~\key{Apps}. &
  Check status. Try reopening the application. \\
Shell déconnecté&
Attendez 60~s pour reconnecter automatiquement.&
  If still disconnected: check Signal~K and Wi-Fi. \\
Profondeur de l'assiette dans l'eau de shoal&
Vérifiez immédiatement avec un instrument de sauvegarde.&
  Slow down. Use a hand lead or backup echo sounder. \\
Le pilote automatique ne répond pas&
  Tap \key{Stand By}. Take manual helm. &
  Investigate the cause before re-engaging. \\
Données inexistantes partout&
Vérifiez l'état de connexion.&
  Allow recovery or check Signal~K server. \\
Données manquantes pour une valeur&
Recoupez-vous avec une autre source.&
  Check the sensor and the Signal~K path. \\
\bottomrule
\end{longtable}

\section*{One-Tap Navigation}

\rowcolors{2}{LtGray}{white}
\begin{longtable}{L{5.8cm}L{7.9cm}}
\toprule
\rowcolor{Navy}
\color{white}\sffamily\bfseries Goal &
\color{white}\sffamily\bfseries Action \\
\midrule
Return to home screen              & Tap \iconlg{applications}~\key{Apps} --- always available \\
Switch to another open app         & Tap its icon in the open-apps strip \\
Change comfort preset              & \ui{Settings~> Comfort} \\
Check connection status            & Status area in Top Bar (right side) \\
Mark anchor position               & Tap \iconlg{anchor-iec}~\key{Anchor} \textrightarrow{} \iconlg{anchor-drop}~\key{Drop} \\
Cancel anchor watch                & Tap \iconlg{anchor-iec}~\key{Anchor} \textrightarrow{} \iconlg{anchor-raise}~\key{Raise} \\
Cancel a POB event                 & POB bar \textrightarrow{} select incident \textrightarrow{} \key{Cancel} \\
Export voyage track                & \ui{My~Data~> Tracks} \textrightarrow{} \key{Export} \\
\bottomrule
\end{longtable}

\chapter{Notes de sécurité}
\label{ch:safety}

\begin{warningbox}
FairWindSK est une aide à la navigation.
Il ne remplace pas le maçonnerie, le veilleur, le papier ou les cartes à moteur indépendant,
ou toute procédure de sécurité certifiée.
Le capitaine et l'équipage du navire restent seuls responsables de la conduite sécuritaire des
le voyage.
\end{warningbox}

\begin{warningbox}
Toujours vérifier les informations essentielles à la sécurité --- position, profondeur, cap,
risque de collision --- avec des instruments indépendants avant d'agir.
Si FairWindSK contredit ce que vous voyez à l'extérieur du bateau ou sur des instruments de confiance,
Faites d'abord confiance à vos yeux et à vos références de navigation indépendantes.
\end{warningbox}

\begin{warningbox}
La profondeur manquante n'est pas la profondeur zéro.
Une valeur manquante signifie que la mesure n'est pas disponible.
Traitez-le comme totalement inconnu et naviguez avec toute la prudence.
\end{warningbox}

\begin{warningbox}
Après une reconnexion ou un redémarrage du serveur, reconfirmez toutes les valeurs de navigation
selon eux.
Donnez au GPS, à la profondeur, au vent et aux capteurs de cap le temps de ré-acquérir et de s'installer.
\end{warningbox}

\begin{warningbox}
Les commandes du pilote automatique doivent être confirmées en fonction du comportement réel du navire.
Ne jamais supposer qu'un changement de mode a pris effet parce que le bouton a répondu.
Tenez-vous prêts à vous désengager et prenez la barre manuelle à tout moment.
\end{warningbox}

\begin{warningbox}
Activer une alarme à FairWindSK soulève une notification Signal~K à bord
réseau.
Il ne contacte ni les garde-côtes, ni les MRCC, ni aucune autorité externe.
Faire un appel de détresse DSC et/ou voix sur VHF Channel~16 en cas d'urgence réelle.
\end{warningbox}

\begin{warningbox}
En cas d'urgence, donner la priorité à la manutention des bateaux, à la sécurité de l'équipage et à l'observation directe
toute interaction avec un écran électronique.
Un écran est un outil; votre savoir-faire est votre sécurité.
\end{warningbox}

\vfill
\begin{center}
  {\color{Teal}\rule{0.5\textwidth}{1pt}}\\[10pt]
  {\Large\sffamily\bfseries\color{Navy}Safe sailing.}\\[8pt]
  {\normalsize\sffamily\color{gray}
    The OpenFairWind Team \textbullet{}
    University of Naples 'Parthenope'}\\[4pt]
  {\small\sffamily\color{Teal}\url{https://github.com/OpenFairWind/FairWindSK}}
\end{center}
